% Standalone problem document, generated from the adjacent README.md.
% Compile with XeLaTeX. Mathematical content is maintained in Markdown.
\documentclass[11pt,a4paper]{article}
\usepackage[fontsize=10.5pt]{scrextend}
\usepackage[margin=22mm,headheight=15pt,headsep=9mm,footskip=11mm]{geometry}
\usepackage{fontspec}
\setmainfont{texgyrepagella}[Extension=.otf,UprightFont=*-regular,BoldFont=*-bold,ItalicFont=*-italic,BoldItalicFont=*-bolditalic]
\setsansfont{texgyreheros}[Extension=.otf,UprightFont=*-regular,BoldFont=*-bold,ItalicFont=*-italic,BoldItalicFont=*-bolditalic]
\setmonofont{texgyrecursor}[Extension=.otf,UprightFont=*-regular,BoldFont=*-bold,ItalicFont=*-italic,BoldItalicFont=*-bolditalic,Scale=MatchLowercase]
\usepackage{amsmath,amssymb}
\usepackage{unicode-math}
\setmathfont{texgyrepagella-math.otf}
\usepackage{xcolor,microtype,enumitem,needspace,fancyhdr}
\usepackage{bookmark}
\definecolor{ink}{HTML}{17344B}
\definecolor{accent}{HTML}{197D80}
\definecolor{muted}{HTML}{596773}
\hypersetup{colorlinks=true,linkcolor=accent,urlcolor=accent,pdftitle={MI-18 - Bapat's
q-permanent monotonicity
conjecture},pdfauthor={Open Problems in Numerical Linear Algebra}}
\urlstyle{same}
\setlength{\parindent}{0pt}
\setlength{\parskip}{5pt plus 1pt minus 1pt}
\linespread{1.04}
\setlength{\emergencystretch}{3em}
\widowpenalty=10000
\clubpenalty=10000
\displaywidowpenalty=10000
\allowdisplaybreaks[1]
\setlist{nosep,leftmargin=1.5em}
\providecommand{\tightlist}{\setlength{\itemsep}{0pt}\setlength{\parskip}{0pt}}
\providecommand{\pandocbounded}[1]{#1}
\pagestyle{fancy}
\fancyhf{}
\fancyhead[L]{\sffamily\footnotesize\color{muted}OPEN PROBLEMS IN NUMERICAL LINEAR ALGEBRA}
\fancyhead[R]{\sffamily\bfseries\footnotesize\color{accent}MI-18}
\fancyfoot[L]{\parbox[b]{0.9\textwidth}{\sffamily\fontsize{8}{10}\selectfont\color{muted}Editorial ratings; a bounded search cannot certify that a problem remains unsolved.\\Literature check: 2026-09-10}}
\fancyfoot[R]{\sffamily\footnotesize\color{muted}\thepage}
\renewcommand{\headrulewidth}{0.3pt}
\renewcommand{\headrule}{\hbox to\headwidth{\color{accent}\leaders\hrule height \headrulewidth\hfill}}
\makeatletter
\renewcommand{\section}{\Needspace{5\baselineskip}\@startsection{section}{1}{0pt}{12pt plus 2pt minus 2pt}{4pt}{\sffamily\bfseries\large\color{ink}}}
\renewcommand{\subsection}{\Needspace{4\baselineskip}\@startsection{subsection}{2}{0pt}{9pt plus 1pt minus 1pt}{3pt}{\sffamily\bfseries\normalsize\color{ink}}}
\makeatother
\setcounter{secnumdepth}{0}
\begin{document}
{\sffamily\small\color{muted}Matrix inequalities and norms\par}
\vspace{3pt}
{\raggedright\hyphenpenalty=10000\exhyphenpenalty=10000\sffamily\bfseries\fontsize{21}{24}\selectfont\color{ink}Bapat's
q-permanent monotonicity conjecture\par}
\vspace{7pt}
{\sffamily\small\textbf{Difficulty:} extreme\quad\textbf{Importance:} interesting
to specialist\par}
{\sffamily\small\bfseries\color{accent}Status: Partially resolved\par}
\vspace{4pt}
{\color{accent}\hrule height 0.8pt}
\vspace{6pt}
\textbf{Rating rationale:} Strict monotonicity on the whole PSD cone is
a longstanding generalized-permanent barrier; its immediate impact is in
specialist matrix-function inequalities.

\subsection{Problem statement}\label{problem-statement}

For a permutation \(\sigma\in S_n\), let \[
\operatorname{inv}(\sigma)=
\#\{(i,j):1\le i<j\le n,\ \sigma(i)>\sigma(j)\}.
\] For real \(q\) and \(A=(a_{ij})\in\mathbb C^{n\times n}\), define \[
P_q(A)=\sum_{\sigma\in S_n}q^{\operatorname{inv}(\sigma)}
\prod_{i=1}^n a_{i,\sigma(i)},\qquad 0^0=1.
\] Is it true that, for every \(n\ge2\), every non-diagonal Hermitian
PSD matrix \(A\) with \(a_{ii}>0\) for all \(i\), and every
\(-1\le q_1<q_2\le1\), \[
P_{q_1}(A)<P_{q_2}(A)?
\] The expression is real for Hermitian \(A\). The order of rows and
columns in the inversion statistic is fixed.

\subsection{Relevance and ratings}\label{relevance-and-ratings}

This polynomial interpolates \(\det A\) at \(q=-1\), the diagonal
product at \(q=0\), and \(\operatorname{per}A\) at \(q=1\). The
conjecture would compare a continuous family of matrix functions on the
PSD cone, extending determinant/permanent inequalities. The non-diagonal
and positive-diagonal hypotheses remove constant-polynomial cases.

\subsection{References}\label{references}

\begin{itemize}
\tightlist
\item
  R. B. Bapat and A. K. Lal, \emph{Inequalities for the q-permanent},
  Linear Algebra and its Applications 197--198 (1994), 397--409,
  monotonicity conjecture
  (\href{https://doi.org/10.1016/0024-3795(94)90497-9}{original paper}).
\item
  L. Mitchell, \emph{A note on Bapat's q-permanent conjecture},
  Operators and Matrices 14 (2020), 915--919, unnumbered Conjecture on
  p.917 and Theorems 1--3
  (\href{https://files.ele-math.com/articles/oam-14-56.pdf}{primary
  paper}).
\item
  C. M. da Fonseca, \emph{The \(\mu\)-permanent revisited} (2018), §4,
  related conjectures and corrections concerning ordered graph cases
  (\href{https://arxiv.org/pdf/1804.02231}{primary manuscript}).
\end{itemize}

\subsection{Status check --- 2026-09-10}\label{status-check-2026-09-10}

Mitchell proves that the original positive-definite conjecture is
equivalent to the displayed extension, and proves rank-one and
order-three cases. Searches for ``Bapat q-permanent conjecture'',
``monotonicity'', ``proof'', ``counterexample'', and 2025--2026 found no
resolution. The stronger proposed extension beyond \([-1,1]\) is not
included; Mitchell explains why extending that assertion to all PSD
matrices fails. No recent explicit reaffirmation of the exact conjecture
was found.

\textbf{Independent audit:} Independently rechecked Mitchell's
PSD-extension conjecture and rank-one theorem, and searched for later
q-permanent monotonicity resolutions. Rank-one and low-order cases are
substantive parts of the displayed family; the general case still has
only historical-source status evidence.
\end{document}
